\scanpage{34}
\mspage{16}
point de vue actuel sont sans importance, puisqu'on travaille dans $\mathrm{Sl}(2,\Lambda)'$) on tombe pile sur les formules du formulaire 1981, qui avaient été obtenues par voie topologico-transcendante (classique) relativement délicate. Cel\unsure{a} confirme que ledit formulaire (où \inserted{le couple} des matrices $\rho=\lambda_\infty^{-1}$ et $\sigma=\lambda_1$ choisies comme générateurs de $\mathrm{Sl}(2,\mathbb Z)$, \deleted{par}
\[
\mathrm{Sl}(2,\mathbb Z)=\{\rho,\sigma\mid\rho^3=\sigma^2=\rho^{-3}\}\tag*{(84)}
\]
\deleted{sont bel et bien}\inserted{n'était défini, \deleted{\unsure{via|par|que par}}} ses propriétés intrinsèques remarquables, qu'à conjugaison près par un élément de $\mathrm{Gl}(2,\mathbb Z)/\{\pm1\}\simeq\mathrm{GP}(1,\mathbb Z)$) \deleted{\unsure{est}}était vraiment ``le plus naturel possible''.

\uline{NB} Si on regarde les homomorphismes
\[
\Pi_{0,3}=\{\rho_0,\rho_1,\rho_\infty\mid\rho_\infty\rho_1\rho_0=1\}\xrightarrow{\varphi^\bullet}\mathrm{Sl}(2,\mathbb Z)\tag*{(85)}
\]
qui relèvent l'homomorphisme précédent à valeurs dans $\mathrm{Sl}(2,\mathbb Z)'$
\[
\rho_0\mapsto\lambda_0^\bullet,\ \rho_1\mapsto\lambda_1^\bullet,\ \rho_\infty\mapsto\lambda_\infty^\bullet,\tag*{(*)}
\]
ou les homomorphismes
\[
\Pi_{0,3}^{\tau}=\{\tau_0,\tau_1,\tau_\infty\mid\tau_{\unsure{0|\infty}}^2=\tau_1^2=\tau_\infty^2\}\xrightarrow{\varphi}\mathrm{Gl}(2,\mathbb Z)\tag*{(86)}
\]
qui relèvent l'hom précédent dans $\mathrm{Gl}(2,\mathbb Z)'\simeq\unsure{\mathrm{GP}(1,\mathbb Z)}$
\[
\tau_0\mapsto r_0^\bullet,\ \tau_1\mapsto r_1^\bullet,\ \tau_\infty\mapsto r_\infty^\bullet\tag*{(**)}
\]
(en partant du fait que pour $(*)$ et $(**)$ il y a des choix canoniques concordants, sans \deleted{\unsure{fixation}}\inserted{\unsure{ambi-|aucune}}
% archive sheet number: 33
